\section{What the principle forces, and what it leaves}\label{sec:essay}

This section is an essay. Its sentences are marked: those about mathematics are statements of the preceding sections,
and those about meaning are readings.

\paragraph{Structure on stated inputs.}
We started from one sentence, that existence is what remains when cancelling to nothing cannot complete, and read it as
a grammar of records. The mathematics of the paper is what that reading forces. A first closed path that mixes the
records is the figure-eight's monodromy, under four stated inputs (Theorem~\ref{thm:seed}). Across all closed paths at
once the moves permute the records' three parities, for every word, and carry a triplet, proved for every odd-trace word, with
the arithmetic of its common point a census to length eight (Theorem~\ref{thm:triplet}). Up to a stated hand and a power of the eta character, the triplet's
symmetry is the multiplier of a classical family of theta functions (Theorem~\ref{thm:theta}). In the fibre frame and
given the triplet postulate the coupling ansatz has a null vector at every point (Theorem~\ref{thm:isotropy}); the coupling's
determinant factors through two defects and their pairing (Theorem~\ref{thm:joint}); and, given three stated inputs, a
reading of the weave's surface and an index formula from the record, the index is $\pm3$ (Proposition~\ref{thm:index}). Nothing in this list was
put in to make a number come out. Its first two statements, the parities and the triplet at the common point, hold for every thread at once (the triplet for every odd-trace one); the later
statements hold in one frame.

\paragraph{Unforced selection.}
What the same mathematics does not force is any of the choices that would single out a world (Section~\ref{sec:selection}).
We located each: no residual symmetry of the weave's group that reduces the count of flavour parameters is compatible with the CKM data at
$3\sigma$; its clock stays in a grading; no frame on record is shown to satisfy the principle's demands on its own; one Higgs at the symmetric point fits the five
non-electron ratios on none of the admissible lines; and the strength of the pull-back sets a side and a speed and no
size. The contrast is the paper's organising one. The structure holds on stated inputs, and it fixes no number. Which structure we are in is not fixed.

\reading{A principle of existence that is minimal cannot also be a principle of selection, because selection would be an
additional content. That the structure holds on its inputs and the ``which'' is not fixed is what one would expect if existence were
the remainder of cancelling, and a particular world were a reading of that remainder. The same boundary would also appear if the chosen form were not the principle's own, since the
frames on record share the same postulated triplet and read-out. This is an interpretation of the
boundary and not a consequence of it.}

\paragraph{A third kind of link.}
The record sorts its identifications by what stands on their right. Those between two structures get decided, one way
or the other. Those from a structure to a value have not been earned. Conditions of consistency are a third kind, and
the moves' anomaly is one of them (Proposition~\ref{prop:anomaly}). If the moves act on the fermions as a symmetry,
they must act on left- and right-handed quarks alike. That constrains the dictionary without identifying anything,
which is also how the Standard Model's own hypercharges are fixed. It does not fix a value (0 of 19).

\reading{The measurer as part of the weave. Under the paper's own rule the measurer is not a thread chosen by hand.
Either every thread measures every other, or the measurer is shared by all of them.
\begin{itemize}
\item \emph{What it needs.} A measurer needs a record that its own motion keeps. No thread keeps a fixed one, and neither does
  the weave as a whole; only a shear or the mirror does (Proposition~\ref{prop:records}). So a measurer whose record
  is a fixed state could sit only at an end, a cusp of the modulus or the mirror. A record carried by a relation
  between changing parts is not excluded by that proposition, and is open. Pointer states are likewise the states that the
  measuring interaction leaves fixed~\cite{Zurek2003}.
\item \emph{What it buys.} For magnitudes read as composing ratios it buys the unit (Lemma~\ref{lem:composition}); what fixes the
  ratios themselves is not addressed. Read this
  way, the overall scale is the measurer's unit, and the side of positivity is the direction in which its records
  accumulate, the reading of time as correlation with a clock inside the system~\cite{PageWootters1983}. Neither input
  is discharged; each changes from a fact about the world to a convention of the measurer.
\item \emph{The other two versions.} Selecting only where a measurer can exist needs a criterion of its own, and the
  record's criterion, a kept fixed record, holds on no thread; selection of this kind gives windows, not
  values~\cite{Weinberg1987,AgrawalBarrDonoghueSeckel1998}. Demanding that the measurer's record of a parameter be the
  parameter needs a kept record; a fixed one exists only at an end, where the record already prices the end data as a
  choice. At the fibre's canonical end the radius is of order one (Census law~\ref{cl:cusp}).
\end{itemize}
This is an interpretation, and it fixes no number.}

\paragraph{The push-back picture, as a question.}
One picture motivated this program from its start, and we record it as a question. In it, every step of the grammar is
accompanied by its pull-back; the two are never at rest; and the interaction of the steps with their pull-backs, and not
either alone, is what is real. The mathematics that bears on the picture is proved and is separate from it:
\begin{itemize}
\item the walk of steps and pull-backs escapes at a positive rate: for equal weights by the theorems of Kesten and
  Furstenberg, and at the other weights by Furstenberg's theorem, with the values measured (Census law~\ref{bd:strength});
\item a full turn round a balance point of a step and its pull-back carries the sign (Section~\ref{sec:grammar});
\item the matter's amplitude along a thread returns by the multiplier of $\eta^6$, so a coherent average survives only on the
  quarter of threads with $\#L\equiv\#R \pmod 4$ (Theorem~\ref{thm:phase});
\item the pull-back's strength sets the probability of the side and the speed and adds no number to the sizes (Census law~\ref{bd:strength});
\item the threads' averages approach the surface's constant mode, and the normalised coherence of the cusp family climbs
  toward one, a read-out (Scoreboard~\ref{bd:layers}).
\end{itemize}
The last two are a sealed negative and a sealed scoreboard, and none is a value (0 of 19). What is not shown is whether the
pull-back is a force, whether the motion is time, whether the values are the motion, and which average the motion takes.

\reading{Is the universe a wave, a weave of threads, or both? In the mathematics the two are one object seen two ways:
the threads are closed geodesics of the modular surface and the waves are its automorphic forms, and the trace formula
equates the sum over the one with the sum over the other. Whether physics lives on this is a question. If it does, the
question to ask next is the one of Section~\ref{sec:selection}: what reads one side through the other, and with what
amplitude.}

\paragraph{What the program stands on.}
It stands on the structure of Sections~\ref{sec:grammar}--\ref{sec:three}, on the boundary of Section~\ref{sec:boundary},
and on its price, every input stated in Table~\ref{tab:inputs}. It does not stand on the readings.
